\scanpage{14}
topologiques il y a un isomorphisme topologique unique qui soit une \deleted{\unsure{isomé}} \inserted{similitude}. (donc une isométrie dans le cas où les surfaces sous-jacentes sont \deleted{\unsure{hémiques|hériques|hyperboliques}} sphériques normalisées ou hyperboliques normalisées)

\inserted{c) Le groupe \unsure{des} autom. métriques d'une carte géod. st. \unsure{$(X,\text{etc.})$|$(X,K,S)$} \inserted{\unsure{métrique} $(X\supset K\supset S,g)$} opère \uline{proprement} sur $X$ -- il \deleted{\unsure{est}} c'est un ss-grpe \uline{discret} de $\operatorname{Aut}(X,g)$}

\uline{NB} On dit que la carte géodésique standard est \uline{normalisée} si a) \unsure{Dans|dans} le cas sphérique ou hyperbolique, $(X,g)$ est normalisé b) dans le cas euclidien -- où, cf plus bas, la carte est \unsure{loc.|loc. f.} finie \deleted{\unsure{sommets et faces,}} donc les longueurs des arêtes sont toutes égales -- on demande que la longueur des arêtes soit égale à $1$ (on pourrait \deleted{\unsure{aussi}} \inserted{à la place} exiger que l'aire des faces soit égale à $1\ldots$). On peut donc \deleted{dire que} \inserted{\uline{Corollaire}} Le foncteur \textquotedblleft oubli de la métrique\textquotedblright\ de la catégorie des cartes \deleted{cellulaires} \inserted{géodésiques} \deleted{\unsure{g}} standard normalisées (avec comme morphismes les isométries qui respectent la carte) dans la catégorie des cartes standard, est une équivalence de catégories.

\uline{Re NB} Carte géodésique standard normalisée \uline{épinglée} : \unsure{Les|Les deux} surfaces ambiantes s'\unsure{identifie-|identifie.|identifient}%
% archive sheet number: 13

